% 원제: 4. 검증 체계
\chapter{검증 체계}

\section{검증 수준과 필요성}

전력변환 시스템의 실패는 서로 다른 시간 규모와 원인에서 나타나므로, 한 종류의 모델로 전체를 판정할 수 없다. 본 과제는 단순한 모델에서 자세한 모델로 올라가도록 검증을 구성했고, 각 수준은 아래 수준이 볼 수 없는 현상을 확인한다(표 \ref{t:levels}). 한 수준의 통과는 그 수준의 가정 안에서만 유효하므로, 위 수준의 완료로 넓혀 해석하지 않는다.

{\footnotesize
\begin{longtable}{L{\dimexpr 0.1746\linewidth-2\tabcolsep\relax}L{\dimexpr 0.2222\linewidth-2\tabcolsep\relax}L{\dimexpr 0.3651\linewidth-2\tabcolsep\relax}L{\dimexpr 0.1270\linewidth-2\tabcolsep\relax}L{\dimexpr 0.1111\linewidth-2\tabcolsep\relax}}
\caption{검증 수준, 필요한 이유와 현재 상태}\label{t:levels}\\
\toprule
\tabhead \textbf{수준} & \textbf{모델·도구} & \textbf{필요한 이유} & \textbf{근거} & \textbf{상태} \\
\midrule
\endfirsthead
\toprule
\tabhead \textbf{수준} & \textbf{모델·도구} & \textbf{필요한 이유} & \textbf{근거} & \textbf{상태} \\
\midrule
\endhead
\bottomrule
\endlastfoot
① 해석·설계 추정 & 평형점, 규격 대조, 데이터시트 손실 규칙 & 정격 분리·보장 하한이 성립해야 제어기 설계가 의미를 갖고, 구조 간 손실 비교는 같은 규칙으로만 공정하다 & E3, E4, E15, E18, E25 & 완료 \\
② 평균모델 & Stage B 6상태 모델, 기준 구조, 창 조정 모듈, 보호 상태기계, 조정형 Stage C, V\textsubscript{A} 노드 & 수천 건의 부하·공차·지연·고장 조합을 빠르게 시험해 제어기·정격·순서를 정한다 & E2, E5–E9, E13, E14, E16, E17, E19–E23, E26 & 완료 \\
③ 독립 구현 & MATLAB 재구현, ngspice & 판정이 한 구현에서만 나오면 코드 오류와 설계 특성을 구별할 수 없다 & E2, E11, E22 & 완료 \\
④ 스위칭 수준 모델 & Stage A 8위상(위상 셰딩 포함), Stage B(5 ns), Stage C 직류변압기(레그 출력용량), DAB 정상상태(데드타임·출력용량) & 리플, 데드타임, DAB 위상 동작, 센서 오차, 영전압 전이처럼 평균모델에 없는 현상을 확인한다 & E1, E10, E12, E16, E21, E22, E24, E26 & 완료 \\
⑤ 소자 수준 회로 & Simulink·Simscape Electrical & 직접 작성한 방정식과 근사를 독립 회로 해석기로 점검하고, 세 단을 한 회로로 연결하며, 결선·게이트 사양을 남긴다 & TB01–TB20 & 완료(대표 사례) \\
⑥ HIL·시제품 & 실시간 시뮬레이터, 축소전력 시제품 & 실제 지연, 잡음, 열, 전자기 간섭 & — & 후속 단계 \\
\end{longtable}
}

\section{판정 원칙}

\begin{itemize}
\item 판정 기준은 시험 전에 고정하고(3장), 결과를 본 뒤에 바꾸지 않는다. 기준을 좁히거나 넓힐 때는 새 시험 번호로 다시 판정한다(예: TB10 회로 기준 2배 → TB11 1.5배). 기준을 아주 조금 넘은 경우도 불통과로 보고한다(예: TB20 사례 D, 1.01 A 대 1.0 A).
\item 모델은 실제 상태를 강제로 자르지 않는다. 제어기는 지령만 제한하므로 과부하나 지연으로 생기는 정격 위반이 그대로 드러난다.
\item 실패한 사례는 지우지 않고 회귀 시험으로 남긴다. 모델 사이의 차이(에지 양자화, 추가 소자, 해석기 특성)는 결과와 함께 명시한다.
\item 소자 수준 회로 모델은 검증 단계마다 새 테스트벤치를 만들어 남기고, 같은 조건의 검증 모델 결과를 기준값으로 둔다. 회로에서 새 현상이 드러나면 먼저 검증 모델로 재현해 원인을 확인한 뒤 설계를 고친다(예: TB14의 영전압 전이, TB17의 Stage A 경부하).
\end{itemize}
